% Abhakseite „Das kann ich“ – Zone nur im Gesamt (\mitzone)
%% TODO Ich-kann-Sätze: Platzhalter wie in den Aufgabendateien
\begin{abhakseite}
\ifdefined\mitzone
\abhakgruppe{Kennst du schon}
\abhak{1}{Möglichkeiten vollständig aufzählen (Liste, Tabelle, Baum ohne Wahrscheinlichkeiten), Zählprinzip „mal“ bei zwei Auswahlen, Anordnungen einer kleinen Menge auflisten und als Produkt zählen, Ergebnismenge als geordnete Paare}
\abhak{2}{Potenzen mit natürlichen Exponenten, Produkte großer Zahlen mit dem Rechner, Ergebnisse in Zehnerpotenzschreibweise lesen und angeben, die Rechnertasten für Fakultät und Binomialkoeffizient (n!, nCr; gerätabhängig)}
\abhak{3}{Brüche und Quotienten kürzen (Fakultätenquotienten, Anteile aus zwei Anzahlen), Anteil als Dezimalzahl und Prozent, Vergleich mit einer Schranke}
\abhak{4}{Pfadregeln}
\abhak{5}{Werte einer Folge mit dem Rechner berechnen und gegen eine Schranke prüfen (beide Nachbarwerte belegen)}
\abhak{6}{Möglichkeiten vollständig aufzählen (Liste, Tabelle, Baum ohne Wahrscheinlichkeiten), Zählprinzip „mal“ bei zwei Auswahlen, Anordnungen einer kleinen Menge auflisten und als Produkt zählen, Ergebnismenge als geordnete Paare – Fehler finden}
\abhak{7}{Möglichkeiten vollständig aufzählen (Liste, Tabelle, Baum ohne Wahrscheinlichkeiten), Zählprinzip „mal“ bei zwei Auswahlen, Anordnungen einer kleinen Menge auflisten und als Produkt zählen, Ergebnismenge als geordnete Paare – selbst rechnen}
\fi
\abhakgruppe{Einheit 1 · Zählprinzip und Anordnungen}
\abhak{8}{Zählprinzip und Anordnungen}
\abhak{9}{Zählprinzip und Anordnungen – Fehler finden, Begründen, Anwendung}
\abhakgruppe{Einheit 2 · Auswahlen und Binomialkoeffizient}
\abhak{10}{Auswahlen und Binomialkoeffizient}
\abhak{11}{Auswahlen und Binomialkoeffizient – Fehler finden, Begründen, Anwendung}
\abhakgruppe{Einheit 3 · Zählen mit Bedingungen und Wahrscheinlichkeitsterme}
\abhak{12}{Zählen mit Bedingungen und Wahrscheinlichkeitsterme}
\abhak{13}{Zählen mit Bedingungen und Wahrscheinlichkeitsterme – Fehler finden, Begründen, Anwendung}
\end{abhakseite}
